\documentclass[11pt, a4paper]{article}

\usepackage{fontspec}
\usepackage{geometry}
\usepackage{fancyhdr}
\usepackage[hidelinks]{hyperref}
\usepackage[normalem]{ulem}
\usepackage{multicol}
\usepackage{enumitem}
\usepackage{amsmath}

\geometry{
  top=2cm,
  bottom=2cm,
  left=2cm,
  right=2cm,
  marginparsep=4pt,
  marginparwidth=1cm
}

\renewcommand{\headrulewidth}{0pt}
\pagestyle{fancyplain}
\fancyhf{}
\lfoot{}
\rfoot{}

\setlength{\parindent}{0pt}
\setlength{\parskip}{0pt}

\usepackage{xunicode}
\defaultfontfeatures{Mapping=tex-text}

\setlist[itemize]{leftmargin=1.4em, itemsep=0pt, topsep=1pt, parsep=0pt}

\newcommand{\code}[1]{\texttt{#1}}

\begin{document}

\small
\sloppy

\begin{center}
\normalsize\textbf{DSST389: Exam III --- Reference Sheet}
\end{center}

\smallskip

\textit{This sheet is attached to your exam. You do not need to memorize any of it.}

\smallskip

\textbf{Building a network.}
A network is a set of \textbf{nodes} and a set of \textbf{edges} joining pairs of them. It is
recorded as an \textbf{edge list}: one row per edge, with the endpoints in the columns
\code{doc\_id} and \code{doc\_id2}. In a \textbf{directed} network the two ends differ; in an
\textbf{undirected} one they do not. Any directed network can be read as undirected by
ignoring the direction; the reverse is not possible.

\begin{verbatim}
node, edge, G = DSNetwork.process(edges, directed=False)
\end{verbatim}

\begin{itemize}
  \item The graph is built \textit{from the edge list}, so a node belonging to no edge is
        never created. Filtering an edge list deletes nodes, silently.
  \item \code{node} has one row per node: \code{id}, the layout coordinates \code{x} and
        \code{y}, and the metrics below. \code{edge} has one row per edge, with \code{x},
        \code{y}, \code{xend}, \code{yend}, which are the four aesthetics of
        \code{geom\_segment}. Add \code{arrow=arrow(length=0.02)} to draw direction.
\end{itemize}

\smallskip

\textbf{Components.}
A \textbf{component} is a set of nodes all reachable from one another by following edges. A
network with exactly one component is \textbf{connected}. The \code{component} column numbers
them in \textbf{descending order of size}, so component 1 is always the largest, and
\code{component\_size} gives the size of each.

\smallskip

\textbf{Centrality.}
All four scores below are computed \textbf{separately within each component}, because the
last three are only defined on a connected network. Filter to \code{component == 1} before
comparing scores across nodes.

\begin{itemize}
  \item \code{degree} --- the number of neighbors a node has. Purely local.
  \item \code{eigen} --- \textit{eigenvalue centrality.} A score proportional to the sum of the
        scores of its neighbors, so that being adjacent to central nodes makes a node central.
        \textbf{Rescaled so the largest score in each component is 1}, which means a score of 1
        says ``the most central node \textit{in this component}'' and cannot be compared across
        components.
        \begin{equation*}
        s_j \;=\; \lambda \cdot \sum_{i \,\in\, \text{Neighbors}(j)} s_i
        \end{equation*}
  \item \code{close} --- \textit{closeness centrality.} Take the shortest-path distance from
        the node to every other node in its component, and add up the \textbf{reciprocals}. A
        neighbor contributes 1, a node four hops away contributes only $1/4$, so the score
        rewards being near \textit{everything}. Undirected networks only.
  \item \code{between} --- \textit{betweenness centrality.} For every pair of other nodes,
        consider the shortest paths between them; count what fraction of them run
        \textit{through} this node, and add that up over all pairs. High for
        \textbf{bridges}: nodes that are not popular but are unavoidable. Degree and
        betweenness can disagree completely about the same node.
  \item \code{cluster} --- a community label, assigned by an algorithm that tries to put more
        edges inside groups than between them. It is an output, not a property: a node on the
        boundary between two groups may be assigned arbitrarily.
\end{itemize}

\smallskip

\textbf{Directed networks.}
Setting \code{directed=True} changes the node table. \code{degree} is replaced by
\code{degree\_out} (links on the page), \code{degree\_in} (links into the page), and
\code{degree\_total}. \code{eigen} and \code{between} are still computed but ignore the
direction. \textbf{\code{close} is not available}, because a distance to a node you cannot
reach is not defined.

\smallskip

\textbf{Kinds of network.}

\begin{itemize}
  \item \textbf{Citation.} An edge whenever one page links to another. Sensitive to page
        length, so summarize by links \textit{in} rather than links \textit{out}.
  \item \textbf{Co-citation.} An edge between two pages whenever some \textit{third} page cites
        \textbf{both}. Always undirected. The \code{count} column holds the number of citing
        pages a pair shares, and edges are usually kept only above a threshold. One page citing
        $n$ items manufactures $n(n-1)/2$ co-citations at a stroke.
  \item \textbf{Distance.} An edge whenever two documents are closer than a cutoff, using the
        angle distances of Chapter 15.
  \item \textbf{Symmetric nearest neighbor.} Take each document's $k$ nearest neighbors, and
        keep an edge only when two documents are \textbf{mutually} in each other's sets. No node
        can then exceed degree $k$, which stops popular documents from becoming hubs.
\end{itemize}

\smallskip

\textbf{Images as arrays.}
\code{cv2.imread(path)} returns a NumPy array of shape \textbf{(height, width, channels)} with
values from 0 to 255. \code{img.mean()} averages over \textbf{every number}, channels included:
a $2 \times 2$ color image has 12 of them, not 4.

\begin{itemize}
  \item \textbf{OpenCV stores channels in BGR order, not RGB.} \code{[255, 0, 0]} is
        \textit{blue} and \code{[0, 0, 255]} is \textit{red}. \code{[0,0,0]} is black and
        \code{[255,255,255]} is white.
  \item \code{cv2.cvtColor(img, cv2.COLOR\_BGR2HSV)} --- convert to hue, saturation, value.
        Same shape. OpenCV scales \textbf{hue to 0--179} (half of degrees, to fit in 8 bits) and
        saturation and value to 0--255.
  \item \code{DSImage.compute\_colors(img\_hsv)} --- the proportion of pixels in each color
        category; the proportions sum to 1. It checks saturation and value \textit{before} hue,
        so that dark and desaturated pixels are called black, white, or gray rather than being
        binned by a hue that means nothing.
  \item \code{DSImage.plot\_image\_grid(df, ncol=10, label\_name="label")} --- display images.
\end{itemize}

\smallskip

\textbf{RGB to HSV.}
Given $R$, $G$, $B$ scaled to $[0, 1]$, let $M = \max(R,G,B)$, $m = \min(R,G,B)$, and let the
chroma be $C = M - m$. Then $V = M$, and:

\begin{equation*}
S = \begin{cases} 0 & \text{if } V = 0 \\[2pt] \dfrac{C}{V} & \text{otherwise} \end{cases}
\qquad\qquad
H = \begin{cases}
0^\circ & \text{if } C = 0 \\[2pt]
60^\circ \times \dfrac{G - B}{C} \bmod 360^\circ & \text{if } M = R \\[4pt]
60^\circ \times \left( \dfrac{B - R}{C} + 2 \right) & \text{if } M = G \\[4pt]
60^\circ \times \left( \dfrac{R - G}{C} + 4 \right) & \text{if } M = B
\end{cases}
\end{equation*}

Hue is an angle on a color wheel: red $0^\circ$, yellow $60^\circ$, green $120^\circ$, cyan
$180^\circ$, blue $240^\circ$, magenta $300^\circ$, wrapping back to red at $360^\circ$. When
$C = 0$ the pixel is gray and \textbf{has no hue}; the formula would divide by zero, and the
$0^\circ$ it returns instead is indistinguishable from red.

\smallskip

\textbf{Images as vectors.}

\begin{itemize}
  \item \code{ViTEmbedder()} --- returns a \textbf{dense} vector of \textbf{768} numbers for an
        image. Unlike a TF-IDF vector, nearly every entry is non-zero and no single dimension
        has a nameable meaning. Only the geometry matters.
  \item \code{umap(df, features=[c.vit]).predict(full=True)} --- project to two dimensions,
        returned as the columns \code{dr0} and \code{dr1}.
  \item \textbf{Transfer learning}: reusing a model trained by someone else, on other images,
        for other labels. What transfers is the \textit{representation}, not the labels.
  \item \code{SigLIPEmbedder()} --- embeds \textbf{images and text into the same space}, with
        \code{.embed\_image(path)} and \code{.embed\_text(string)}. The vectors are unit length,
        so \code{dot\_product(a, b)} is exactly the cosine similarity of Chapter 15: 1 means
        identical, 0 means orthogonal, and the angle distance is $\arccos$ of it.
  \item \textbf{Zero-shot classification}: embed each category description as text, and give
        each image the label whose text vector it has the largest dot product with. No training,
        no labeled examples. It is an \textit{argmax over the candidate list you supplied}, so
        there is no ``none of the above,'' and adding a category can change a prediction that
        was already made.
\end{itemize}

\end{document}
